\newcommand{\nomeEstado}{BAHIA}
\newcommand{\siglaEstado}{BA}
\newcommand{\nomeConcessionaria}{NEOENERGIA BAHIA}
\newcommand{\cnpjConcessionaria}{15.139.629/0001-94}
\newcommand{\termoCIP}{Sexto}
\newcommand{\presidenteConcessionaria}{Sr. Thiago Guth.}

\newcommand{\empresaResponsavel}{ABEL}
\newcommand{\nomeMunicipio}{Ituberá}
\newcommand{\nomeMunicipioMaisculo}{\MakeUppercase{\nomeMunicipio}}
\newcommand{\IPestimada}{2641033}
\newcommand{\cnpjMunicipio}{14.195.333/0001-28}
\newcommand{\telefone}{(73) 3256-8100}
\newcommand{\email}{administracao@itubera.ba.gov.br}
\newcommand{\sedePrefeitura}{Rua Coronel Barachísio Lisboa, nº 91, Centro, Ituberá – BA, CEP 45435-000}
\newcommand{\nomePrefeito}{Reges Jonas Aragão Santos}
\newcommand{\generoPrefeitura}{pelo Prefeito}
\newcommand{\numeroContrato}{090/2025}
\newcommand{\quantidadeAditivos}{1}
\newcommand{\legitimidade}{1}
\newcommand{\publicacao}{Diário Oficial do Município}
\newcommand{\dataPublicacao}{29 de Abril de 2025}
\newcommand{\tituloClausula}{CLÁUSULA SEXTA - DAS OBRIGAÇÕES DA CONTRATADA}
\newcommand{\clausula}{6.3 — Deverá a empresa efetuar levantamento dos créditos a que faz jus o Município, referentes aos pagamentos indevidos à concessionária de energia elétrica, conforme os prazos estabelecidos na Resolução Normativa da ANEEL nº 1.000, de 7 de dezembro de 2021, especialmente o art. 323, § 2º, inciso I, e suas alterações posteriores, podendo, para tanto, representar administrativamente o Município de Ituberá/BA perante a Agência Nacional de Energia Elétrica — ANEEL, perante a distribuidora de energia elétrica e, caso necessário, perante a agência reguladora conveniada do Estado, em todos os processos, reclamações, recursos e demais atos decorrentes de suas competências.}
